\chapter{Simple Regression}\label{ch:b1:simple-regression}

Among the 2,204 men of the PSID sample with positive earnings in 1978, each
extra year of schooling goes with log earnings higher by $0.0885$: about $9\%$
more earnings. % figdata: fit
Does a year of school raise earnings by $9\%$? The line through the data does
not say. This chapter builds the line --- least squares --- and shows how to
read it, how precise it is, when it estimates what we want, and two
reasons, \omterm{def:b1:simple-regression:toward-mean}{regression toward the mean} and \omterm{def:b1:simple-regression:omitted}{omitted variables}, why it often
does not.

\begin{openingchart}{Mean log earnings in 1978 by years of schooling, PSID comparison
sample, men with positive earnings and 8 to 17 years of schooling (counts
above the points), and the least-squares line fitted on all 2,204 men.}{National
Supported Work comparison sample (PSID), as assembled by LaLonde (1986) and Dehejia and Wahba (1999)}
\begin{tikzpicture}
\begin{axis}[omscatter, xlabel={Years of schooling}, ylabel={Log earnings}, xmin=7.5, xmax=17.5, ymin=9.3,
  ymax=10.5, xtick={8,...,17}, width=0.66\linewidth, height=0.38\linewidth]
  \addplot[omThm, thick, sharp plot, mark=none] table[x=education, y=line, col sep=comma] {figdata/bachelor-1/32-simple-regression/out/school.csv};
  \addplot[only marks, mark=*, mark size=2pt, omDef, nodes near coords, point meta=explicit symbolic,
    nodes near coords style={font=\tiny, black, above=1pt}] table[x=education, y=mean, meta=count, col sep=comma] {figdata/bachelor-1/32-simple-regression/out/school.csv};
\end{axis}
\end{tikzpicture}
\end{openingchart}

The means rise with schooling and lie close to the line; the line explains
$13\%$ of the variance of log earnings across men, and the rest is everything
else. % figdata: fit

\begin{recall}
Covariance, correlation and \omterm{def:b1:describing-data:conditional-mean}{conditional means} are \cref{ch:b1:describing-data};
\omterm{def:b1:random-variables-for-econometrics:conditional-expectation}{conditional expectation} and its best-predictor property,
\cref{ch:b1:random-variables-for-econometrics}; \omterm{def:b1:sampling-and-inference:standard-error}{standard errors} and tests, \cref{ch:b1:sampling-and-inference}.
That a correlation is not a cause was met at school.
\end{recall}

\begin{measure}[Earnings in the PSID and CPS samples]
The two samples are comparison groups drawn from two large US surveys and
assembled to evaluate a job-training experiment (LaLonde 1986; Dehejia and
Wahba 1999): men's earnings in 1974, 1975 and 1978, in dollars of 1982, with
age, schooling and a few other traits. % ledger: b1:lalonde1986
% ledger: b1:nsw-1982dollars
% ledger: b1:dehejia1999
The CPS sample's earnings are capped (\cref{ch:b1:describing-data}); the PSID
sample's are not.
\end{measure}

\section{The conditional expectation function and its best line}\label{sec:b1:simple-regression:population}

\begin{definition}[Population regression line]\label{def:b1:simple-regression:population-line}
\index{population regression line}\index{best linear predictor}
The \emph{population regression line}, or \emph{best linear predictor}, of $Y$ given $X$ is
the line $\beta_0 + \beta_1 X$ that minimises $\E[(Y - b_0 - b_1X)^2]$ over $b_0$ and $b_1$.
\end{definition}

\begin{proposition}[Its coefficients]\label{prop:b1:simple-regression:population}
$\beta_1 = \Cov(X, Y)/\Var(X)$ and $\beta_0 = \E[Y] - \beta_1\E[X]$.
\end{proposition}

\begin{proof}
Setting the derivatives with respect to $b_0$ and $b_1$ to zero gives $\E[Y - b_0 - b_1X] = 0$
and $\E[X(Y - b_0 - b_1X)] = 0$. The first gives $\beta_0$. Substituting it into the second,
$\E[X(Y - \E Y)] - b_1\E[X(X - \E X)] = 0$, that is $\Cov(X, Y) = b_1\Var(X)$. The function is convex in
$(b_0, b_1)$, so this is the minimum.
\end{proof}

The \omterm{def:b1:random-variables-for-econometrics:conditional-expectation}{conditional expectation} function is the best predictor of all
(\cref{prop:b1:random-variables-for-econometrics:best}); the regression line is the best straight line. When the
\omterm{def:b1:random-variables-for-econometrics:conditional-expectation}{conditional expectation} function is a line, the two coincide; otherwise
the line is the best linear approximation to it, as in the opening chart.

\section{Least squares}\label{sec:b1:simple-regression:ols}

\begin{definition}[Ordinary least squares estimator]\label{def:b1:simple-regression:ols}
\index{ordinary least squares}\index{OLS}
The \emph{ordinary least squares} (OLS) \omterm{def:b1:sampling-and-inference:estimator}{estimator} of the line in a sample $(x_i, y_i)$ minimises
$\sum_i (y_i - b_0 - b_1x_i)^2$; it is
\[
  \hat\beta_1 = \frac{\sum_i (x_i - \bar x)(y_i - \bar y)}{\sum_i (x_i - \bar x)^2}, \qquad \hat\beta_0 = \bar y - \hat\beta_1\bar x,
\]
the sample analogue of \cref{prop:b1:simple-regression:population}.
\end{definition}

\begin{definition}[Fitted value, residual]\label{def:b1:simple-regression:residual}
\index{fitted value}\index{residual}
The \emph{fitted value} of observation $i$ is $\hat y_i = \hat\beta_0 + \hat\beta_1x_i$; its \emph{residual} is
$\hat u_i = y_i - \hat y_i$.
\end{definition}

\begin{proposition}[Residuals]\label{prop:b1:simple-regression:residuals}
$\sum_i \hat u_i = 0$ and $\sum_i x_i\hat u_i = 0$: the \omterm{def:b1:simple-regression:residual}{residuals} have mean zero and are
uncorrelated with $x$ in the sample.
\end{proposition}

\begin{proof}
These are the two first-order conditions of the minimisation, $\sum_i (y_i - \hat\beta_0 - \hat\beta_1x_i) = 0$
and $\sum_i x_i(y_i - \hat\beta_0 - \hat\beta_1x_i) = 0$. Since the \omterm{def:b1:simple-regression:residual}{residuals} sum to zero, their sample
covariance with $x$ is $\frac1{n-1}\sum_i x_i\hat u_i = 0$.
\end{proof}

\begin{omfigure}
\begin{tikzpicture}[scale=0.9, font=\small]
  \draw[->] (0,0) -- (7,0) node[right] {$x$};
  \draw[->] (0,0) -- (0,4.4) node[above] {$y$};
  \draw[omThm, thick] (0.3,0.9) -- (6.7,3.8) node[right] {$\hat y = \hat\beta_0 + \hat\beta_1 x$};
  \foreach \x/\y in {1/1.9, 2/1.2, 3/2.6, 4/2.0, 5/3.5, 6/2.9} {
    \pgfmathsetmacro{\f}{0.9 + (\x - 0.3)*2.9/6.4}
    \draw[gray, dashed] (\x,\y) -- (\x,\f);
    \fill[omDef] (\x,\y) circle (2pt);
  }
  \node[gray, font=\scriptsize, right] at (5.05,3.0) {$\hat u_i$};
\end{tikzpicture}

{\small Least squares chooses the line that makes the sum of the squared
vertical distances (the \omterm{def:b1:simple-regression:residual}{residuals}, dashed) as small as possible.}
\end{omfigure}

\omcode{code/bachelor-1/32-simple-regression/listings/ols_by_hand.py}{1}{13}{Least squares from
sums.}\label{lst:b1:simple-regression:by-hand}

\omoutput{code/bachelor-1/32-simple-regression/outputs/ols_by_hand.txt}{The slope, and the two
properties of the \omterm{def:b1:simple-regression:residual}{residuals} up to rounding error.}\label{out:b1:simple-regression:by-hand}

\omcode{code/bachelor-1/32-simple-regression/listings/ols_statsmodels.py}{1}{10}{The same regression
with statsmodels.}\label{lst:b1:simple-regression:statsmodels}

\omoutput{code/bachelor-1/32-simple-regression/outputs/ols_statsmodels.txt}{Coefficients,
\omterm{def:b1:sampling-and-inference:standard-error}{standard errors}, $t$ statistics, $p$-values and $95\%$ intervals.}\label{out:b1:simple-regression:statsmodels}

\section{Reading a slope}\label{sec:b1:simple-regression:reading}

The slope's meaning depends on whether each variable enters in levels or
in logarithms. With $\log y$ and $x$ in levels, one more unit of $x$ goes with
$100\beta_1$ per cent more $y$ (approximately, for small $\beta_1$; exactly $100(e^{\beta_1} - 1)$). With
$y$ and $\log x$, $1\%$ more $x$ goes with $\beta_1/100$ more units of $y$. With both in logs, $\beta_1$ is
an elasticity: $1\%$ more $x$ goes with $\beta_1$ per cent more $y$.

\omcode{code/bachelor-1/32-simple-regression/listings/reading_slopes.py}{1}{16}{Four functional forms,
one pair of variables.}\label{lst:b1:simple-regression:reading}

\omoutput{code/bachelor-1/32-simple-regression/outputs/reading_slopes.txt}{1,849 dollars a year
of schooling; $8.75\%$; 168 dollars per $1\%$; an elasticity of $0.83$ (men with at least
one year of schooling).}\label{out:b1:simple-regression:reading}

\section{Fit}\label{sec:b1:simple-regression:fit}

\begin{definition}[Coefficient of determination]\label{def:b1:simple-regression:r2}
\index{coefficient of determination}
The \emph{coefficient of determination} is $R^2 = \sum_i (\hat y_i - \bar y)^2 / \sum_i (y_i - \bar y)^2 = 1 - \sum_i \hat u_i^2/\sum_i (y_i - \bar y)^2$: the
share of the sample variance of $y$ that the \omterm{def:b1:simple-regression:residual}{fitted values} reproduce. In a simple
regression it equals the squared correlation of $x$ and $y$.
\end{definition}

\omcode{code/bachelor-1/32-simple-regression/listings/r_squared.py}{1}{15}{$R^2$ three
ways.}\label{lst:b1:simple-regression:r2}

\omoutput{code/bachelor-1/32-simple-regression/outputs/r_squared.txt}{The same number three
times.}\label{out:b1:simple-regression:r2}

A low $R^2$ says that $x$ leaves much of the variation in $y$ unexplained; it does
not say that the slope is small, badly estimated or unimportant. Schooling
explains $13\%$ of the variance of log earnings, and the slope is still
measured with a $t$ statistic of 18.

\section{When is the slope unbiased?}\label{sec:b1:simple-regression:unbiased}

\begin{definition}[Error term]\label{def:b1:simple-regression:error}
\index{error term}
In the model $Y = \beta_0 + \beta_1X + u$, the \emph{error term} $u$ collects everything other than $X$
that affects $Y$. For observation $i$ it is $u_i$, unobserved; the \omterm{def:b1:simple-regression:residual}{residual} $\hat u_i$ is
its estimate.
\end{definition}

\begin{definition}[Zero conditional mean assumption]\label{def:b1:simple-regression:zcm}
\index{zero conditional mean assumption}
The \emph{zero \omterm{def:b1:describing-data:conditional-mean}{conditional mean} assumption} is $\E[u \mid X] = 0$: whatever the value of $X$,
the other influences on $Y$ average to zero.
\end{definition}

\begin{model}{The simple linear regression model}
\index{simple linear regression model}
\assumes $Y = \beta_0 + \beta_1X + u$ with $\E[u \mid X] = 0$; a \omterm{def:b1:random-variables-for-econometrics:random-sample}{random sample} $(x_i, y_i)$; some
variation in $x$.
\predicts $\hat\beta_1$ is unbiased for $\beta_1$, with the \omterm{def:b1:sampling-and-inference:standard-error}{standard error} of
\cref{prop:b1:simple-regression:se} when the variance of $u$ does not depend on $X$.
\tested Consistent with the published estimates: the log-earnings slope
of $0.0885$ is close to those of studies that do not rest on the
zero-conditional-mean assumption (the record \emph{Returns to schooling}
below). % ledger: b1:card1999
\breaks When ability or family background are in $u$ and also move schooling
(the third university year returns to schooling with instruments), or
when $X$ is measured with error (second university year).
\end{model}

\begin{proposition}[Unbiasedness]\label{prop:b1:simple-regression:unbiased}
Under the model's assumptions, $\E[\hat\beta_1 \mid x_1, \dots, x_n] = \beta_1$.
\end{proposition}

\begin{proof}
Substituting $y_i - \bar y = \beta_1(x_i - \bar x) + (u_i - \bar u)$ into $\hat\beta_1$ gives
$\hat\beta_1 = \beta_1 + \sum_i (x_i - \bar x)u_i / \sum_i (x_i - \bar x)^2$, since $\sum_i (x_i - \bar x)\bar u = 0$. Given the $x_i$, the
weights are numbers and $\E[u_i \mid x_1, \dots, x_n] = 0$ by \omterm{def:b1:random-variables-for-econometrics:random-sample}{random sampling} and the \omterm{def:b1:simple-regression:zcm}{zero
conditional mean assumption}, so the second term has expectation zero.
\end{proof}

\begin{proposition}[Standard error of the slope]\label{prop:b1:simple-regression:se}
If moreover $\Var(u \mid X) = \sigma^2$, then $\Var(\hat\beta_1 \mid x) = \sigma^2/\sum_i (x_i - \bar x)^2$, and the \omterm{def:b1:sampling-and-inference:standard-error}{standard error} is
estimated by $\sqrt{\hat\sigma^2/\sum_i (x_i - \bar x)^2}$ with $\hat\sigma^2 = \sum_i \hat u_i^2/(n - 2)$.
\end{proposition}

\admitted{The variance follows from the expression for $\hat\beta_1$ in the previous proof; the
$n - 2$ and the case of a variance that changes with $X$ are derived in the second
university year.}

\omcode{code/bachelor-1/32-simple-regression/listings/slope_se.py}{1}{14}{The formula against
statsmodels, and two $t$ tests.}\label{lst:b1:simple-regression:se}

\omoutput{code/bachelor-1/32-simple-regression/outputs/slope_se.txt}{The slope is far from zero,
and also significantly below $0.10$.}\label{out:b1:simple-regression:se}

\section{Regression toward the mean}\label{sec:b1:simple-regression:toward}

\begin{definition}[Regression toward the mean]\label{def:b1:simple-regression:toward-mean}
\index{regression toward the mean}
\emph{Regression toward the mean} is the tendency of units with extreme values of a
variable at one date to have values closer to the mean at another, whenever
the two are imperfectly correlated: with equal variances, the slope of one
on the other is $r < 1$.
\end{definition}

\omcode{code/bachelor-1/32-simple-regression/listings/toward_mean.py}{1}{16}{Earnings in 1978 on
1975, and 1975 on 1978.}\label{lst:b1:simple-regression:toward}

\omoutput{code/bachelor-1/32-simple-regression/outputs/toward_mean.txt}{Both slopes below one;
their product is $r^2$.}\label{out:b1:simple-regression:toward}

\begin{omchart}{Mean earnings in 1978 against mean earnings in 1975, by bins of 1975
earnings (5,000 dollars wide, bins of 20 men or more), PSID comparison sample, and the 45-degree line
(dashed).}{National Supported Work comparison sample (PSID)}
\begin{tikzpicture}
\begin{axis}[omscatter, xlabel={Earnings in 1975 (dollars of 1982)}, ylabel={Mean earnings in 1978 (dollars of 1982)}, xmin=0,
  xmax=55000, ymin=0, ymax=55000, width=0.56\linewidth, height=0.46\linewidth, scaled ticks=false]
  \addplot[black, dashed, sharp plot, mark=none, domain=0:55000] {x};
  \addplot[only marks, mark=*, mark size=2pt, omDef] table[x=re75, y=re78, col sep=comma] {figdata/bachelor-1/32-simple-regression/out/toward.csv};
\end{axis}
\end{tikzpicture}
\end{omchart}

Average earnings rose between the two years, yet the gains were largest at
the bottom, and in the highest bins they mostly turned into losses (three of
the five bins above 30,000 dollars lost): the men of the lowest bin went
from 946 to 6,525 dollars on average, those of the highest from 52,086 to
47,719. % figdata: toward
Nothing pulled them together. A man's earnings in one year combine a lasting
part and a transitory part; men selected for high 1975 earnings include many
with a lucky year, which does not repeat. Reverse the regression and the
same happens backwards: the slope of 1975 on 1978 is also below one.

\begin{record}[The regression fallacy]
In 1933 a book claimed that business enterprises were converging towards
the average: groups of firms ranked by their first-year value moved towards the
average. A review answered that ranked by their last-year value the lines
would diverge, and that the real test of convergence is a consistent
decline in the variance among individual enterprises. Friedman (1992)
quotes the review and finds the same fallacy in comparisons of countries'
growth sixty years later. % ledger: b1:friedman1992
\end{record}

\section{A regression is not a cause}\label{sec:b1:simple-regression:cause}

\begin{definition}[Omitted variable, confounder]\label{def:b1:simple-regression:omitted}
\index{omitted variable}\index{confounder}
An \emph{omitted variable}, or \emph{confounder}, is a variable left out of the regression
that affects $Y$ and is correlated with $X$. It sits in $u$ and breaks the \omterm{def:b1:simple-regression:zcm}{zero
conditional mean assumption}.
\end{definition}

\begin{center}
\begin{tikzpicture}[font=\small, every node/.style={draw, rounded corners, minimum width=2.2cm, minimum height=0.8cm}]
  \node (a) at (2.5,1.8) {Ability};
  \node (s) at (0,0) {Schooling};
  \node (y) at (5,0) {Earnings};
  \draw[->, thick] (a) -- (s);
  \draw[->, thick] (a) -- (y);
  \draw[->, thick, omThm] (s) -- node[draw=none, below] {the effect sought} (y);
\end{tikzpicture}
\end{center}

\omcode{code/bachelor-1/32-simple-regression/listings/omitted.py}{1}{17}{Ability left out: a
simulation with a known answer.}\label{lst:b1:simple-regression:omitted}

\omoutput{code/bachelor-1/32-simple-regression/outputs/omitted.txt}{A true effect of $0.060$
estimated as $0.084$ on average.}\label{out:b1:simple-regression:omitted}

\begin{omfigure}
\begin{tikzpicture}
\begin{axis}[omrecord, xlabel={Estimated slope}, ylabel={Density}, xmin=0.03, xmax=0.11, ymin=0, ymax=62,
  width=0.6\linewidth, height=0.34\linewidth, scaled x ticks=false,
  xticklabel style={/pgf/number format/.cd, fixed, precision=2},
  legend style={font=\scriptsize, at={(0.98,0.96)}, anchor=north east}]
  \addplot[omDef, thick, const plot mark mid] table[x=mid, y=long, col sep=comma] {figdata/bachelor-1/32-simple-regression/out/omitted.csv};
  \addplot[omThm, thick, const plot mark mid] table[x=mid, y=short, col sep=comma] {figdata/bachelor-1/32-simple-regression/out/omitted.csv};
  \draw[black, dashed] (axis cs:0.06,0) -- (axis cs:0.06,58);
  \legend{ability included, ability left out}
\end{axis}
\end{tikzpicture}

{\small Slopes from 1,000 simulated samples of 500 people: with ability in the
regression they centre on the true $0.060$ (dashed); left out, on $0.084$.
% figdata: omitted_means
}
\end{omfigure}

The bias is the \omterm{def:b1:simple-regression:omitted}{omitted variable}'s effect on $Y$ times the slope of the \omterm{def:b1:simple-regression:omitted}{omitted
variable} on $X$: $0.10 \times 1.5/(1.5^2 + 2^2) = 0.024$. The second university year derives this
formula in general.

\begin{definition}[Reverse causation]\label{def:b1:simple-regression:reverse}
\index{reverse causation}
There is \emph{reverse causation} when $Y$ affects $X$, so that a slope of $Y$ on $X$ mixes
the effect of $X$ on $Y$ with the effect of $Y$ on $X$.
\end{definition}

\omcode{code/bachelor-1/32-simple-regression/listings/reverse.py}{1}{12}{Prices and quantities
moved by both curves.}\label{lst:b1:simple-regression:reverse}

\omoutput{code/bachelor-1/32-simple-regression/outputs/reverse.txt}{A slope of $-0.6$: neither the
demand curve's $-1$ nor the supply curve's $+1$.}\label{out:b1:simple-regression:reverse}

Price and quantity are set together by supply and demand. Regressing one on
the other gives a mixture of the two curves, weighted by which one the
shocks moved more. In this market the supply shocks have twice the standard
deviation of the demand shocks, so the points trace out mainly the demand
curve, and the slope is $(1 - 4)/(1 + 4) = -0.6$.

\omcode{code/bachelor-1/32-simple-regression/listings/cps_cap.py}{1}{22}{The same regression on
the capped sample, and the PSID sample capped at the same level.}\label{lst:b1:simple-regression:cps}

\omoutput{code/bachelor-1/32-simple-regression/outputs/cps_cap.txt}{Capping flattens the line:
the PSID slope falls from $0.0885$ to $0.0621$ when its earnings are capped at
the CPS maximum.}\label{out:b1:simple-regression:cps}

Data problems bias slopes too. In the CPS sample $30\%$ of the men with
earnings sit at the maximum the survey records, and its slope, $0.0433$, is
half the PSID's. But the two samples are different populations, so that
comparison mixes the cap with everything else that differs between them.
The cap's own effect is measured within one sample: capping the PSID men at
the CPS maximum, which binds for $40.7\%$ of them, lowers the slope from
$0.0885$ to $0.0621$. Capping the highest earnings, which are those of the
most schooled, flattens the line. % computed: code/bachelor-1/32-simple-regression/outputs/cps_cap.txt

\begin{record}[Returns to schooling]
A regression of log earnings on schooling across people gives slopes like
the opening chart's. Card (1999) reviews the attempts to estimate the
causal return: studies that use changes in the schooling system as
natural experiments find returns $20$--$40\%$ above the corresponding \omterm{def:b1:simple-regression:ols}{OLS}
estimates, and studies of twins suggest an upward ability bias of about
$10\%$. % ledger: b1:card1999
The cross-sectional slope is not the causal effect, but for the standard
samples of these studies it is not far from it. The next volumes treat experiments, instruments and
panels, the tools behind those estimates.
\end{record}

\section*{Exercises}\label{sec:b1:simple-regression:exercises}

\begin{exercise}[$\star$]\label{exo:b1:simple-regression:five}
Find the \omterm{def:b1:simple-regression:ols}{OLS} slope and intercept for the points $(8, 9.4)$, $(10, 9.6)$, $(12, 9.9)$,
$(14, 10.0)$, $(16, 10.3)$.
\end{exercise}

\begin{exercise}[$\star$]\label{exo:b1:simple-regression:read}
From \cref{out:b1:simple-regression:statsmodels}, give the slope, its \omterm{def:b1:sampling-and-inference:standard-error}{standard error} and $R^2$.
\end{exercise}

\begin{exercise}[$\star$]\label{exo:b1:simple-regression:interpret}
A regression of log wages on years of experience has a slope of $0.06$.
Interpret it.
\end{exercise}

\begin{exercise}[$\star$]\label{exo:b1:simple-regression:predict}
Predict the fitted log earnings of a PSID man with 16 years of schooling, and
the earnings it corresponds to.
\end{exercise}

\begin{exercise}[$\star\star$]\label{exo:b1:simple-regression:residuals}
Prove that the \omterm{def:b1:simple-regression:residual}{residuals} of an \omterm{def:b1:simple-regression:ols}{OLS} regression with an intercept are
uncorrelated with the \omterm{def:b1:simple-regression:residual}{fitted values}.
\end{exercise}

\begin{exercise}[$\star\star$]\label{exo:b1:simple-regression:ages}
Run the regression of log earnings on schooling for PSID men under 30 and for
men over 40. Compare the slopes, with their \omterm{def:b1:sampling-and-inference:standard-error}{standard errors}.
\end{exercise}

\begin{exercise}[$\star\star$]\label{exo:b1:simple-regression:energy}
Estimate the income elasticity of energy use in 2019 by a log--log regression
on the countries of \cref{ch:b1:comparative-statics}, with its \omterm{def:b1:sampling-and-inference:standard-error}{standard error}.
\end{exercise}

\begin{exercise}[$\star\star$]\label{exo:b1:simple-regression:direction}
Family wealth raises both children's schooling and their later earnings,
through contacts. Predict the direction of the bias in the slope of earnings
on schooling.
\end{exercise}

\begin{exercise}[$\star\star$]\label{exo:b1:simple-regression:test}
Test at $5\%$ whether the PSID slope equals $0.10$.
\end{exercise}

\begin{exercise}[$\star\star\star$]\label{exo:b1:simple-regression:derive}
Derive the \omterm{def:b1:simple-regression:ols}{OLS} \omterm{def:b1:sampling-and-inference:estimator}{estimator} by minimising the sum of squared \omterm{def:b1:simple-regression:residual}{residuals}, and
prove \cref{prop:b1:simple-regression:unbiased} in your own steps.
\end{exercise}

\begin{exercise}[$\star\star\star$]\label{exo:b1:simple-regression:product}
Show that the slope of $y$ on $x$ times the slope of $x$ on $y$ equals $r^2$, and use it to
explain why both slopes in \cref{out:b1:simple-regression:toward} are below one.
\end{exercise}

\begin{exercise}[$\star\star\star$]\label{exo:b1:simple-regression:refute}
What evidence would refute reading the $0.0885$ slope as the effect of a year
of school?
\end{exercise}

\section{Problem: Schooling and Earnings in Brenland}\label{sec:b1:simple-regression:problem}

\begin{problem}[{Weekend problem --- the Brenland return to a year of school}]\label{pb:b1:simple-regression:brenland}
Brenland's household survey (invented) records the head's years of
schooling and the household's gross income. Use the households with
positive gross income. Part III uses a simulated cohort with an ability
score, written by the same script with a fixed seed: true return $0.05$,
ability effect $0.08$.

\medskip\noindent\textbf{Part I --- The regression.}
\begin{enumerate}
  \item How many households have positive gross income?
  \item Give the mean and standard deviation of schooling.
  \item Regress log gross income on schooling: give the slope and intercept.
  \item Check that the \omterm{def:b1:simple-regression:residual}{residuals} sum to zero.
  \item Give the fitted log income at 12 and at 16 years.
  \item By what factor does fitted income rise from 12 to 16 years?
  \item Give the slope of income in dalens on schooling.
\end{enumerate}

\medskip\noindent\textbf{Part II --- Reading, fit and precision.}
\begin{enumerate}[resume]
  \item Interpret the slope of item 3.
  \item Give the exact percentage it implies.
  \item Give $R^2$.
  \item Does a low $R^2$ mean that schooling matters little for income?
  \item Give the \omterm{def:b1:sampling-and-inference:standard-error}{standard error} and the $t$ statistic of the slope.
  \item Give a $95\%$ interval.
  \item Test whether the slope equals $0.10$.
\end{enumerate}

\medskip\noindent\textbf{Part III --- Hidden ability (simulated cohort).}
\begin{enumerate}[resume]
  \item Give the slope of ability on schooling.
  \item Give the slope of log income on schooling alone.
  \item Give the slopes of schooling and ability when both enter.
  \item Check that the short slope minus the long one equals the ability
    coefficient times the slope of item 15.
  \item By how much does the short slope overstate the true $0.05$?
  \item Why is the bias upward?
  \item Could the survey's slope of item 3 be biased in the same way?
\end{enumerate}

\medskip\noindent\textbf{Part IV --- A second variable (reported, not solved).}
\begin{enumerate}[resume]
  \item Add the head's age to the regression: give the new schooling slope.
  \item Give the age coefficient.
  \item Give the correlation of schooling and age, and say why the slope
    barely moves.
  \item When would adding a variable change the slope a lot?
  \item State the result.
\end{enumerate}
\end{problem}
